% Figure polypeps_cell_statement.
% Owner: none (note-only figure).
% Used by: gapped-2d-notes PP Sec. 2 (Lemma 2, the statement; step 1 cites the right
% panel for the exterior colouring).
% Source: arXiv OpenAI2026PolynomialPEPS, sections/02-information.tex:157-176
% (lem:cell-cmi: G one cell of an axis-parallel grid of side s, B any union of
% other cells, C in Lambda \ (B u G) with dist_inf(C, closed ambient G) >= cs) and
% 179-183 (the exterior colouring of the proof).
% Ink-to-index: dots are lattice sites; there are no tensors or open indices.
% Both panels draw the same 8 x 8 sites (rows 2-9; rows 1 and 10 are blank label
% rows with inkless anchors only), cut into a 4 x 4 grid of cells of 2 x 2 sites.
% G is cell (2,2) of the grid, sites (4,3)..(5,4). Left panel (the statement): G
% (red tint, the set the lemma bounds, label G east of it); B (teal tint, one
% enclosure, label B under the panel) is the union of the cells (1,2), (2,1) and
% (4,1) and of (3,3), which touches G only at a corner; C (blue tint, label C
% above it) is the cells (1,4) and (2,4); the top bracket, from G's east column
% to C's west column, is the clearance dist_inf(C, closed ambient G) >= cs.
% Right panel (step 1 of the proof): every cell of B is forbidden 0 (teal tint
% on top), every other exterior cell, C included, is forbidden 1 (olive tint:
% the whole field minus G, drawn whole with B on top, so no moat); C keeps a
% blue contour, unlabelled (L4); G is the untinted hole, the allowed region of
% the proof. Since the cells are aligned and of equal side, each side of G
% borders exactly one exterior cell, so the exterior colour along the boundary
% of G is constant on each side and can change only at the four corners of the
% ambient cell (red centre dots): here north and west are forbidden 0, east and
% south forbidden 1, and the corner cell (3,3) is forbidden 0. These corners
% are the chamber-line intersections where step 2 reserves its holes.
% Check opt-out: the joiner \longrightarrow is a step of the proof (colour the
% exterior), not a relation; both panels have empty boundary signature.
% Boundary signature: empty in both panels.
% Not to scale: s >= 1 is any side and c > 0 is fixed; a cell of side s is
% drawn as 2 x 2 sites (the plan's 3 x 3 cells would put two 12-site panels at
% 3 mm, below the 4.5 mm lattice floor of shared-spec 4.1 P1); B and C are
% arbitrary (C need not be a union of cells).
% Label macros: none.
% Fit: one column (\tnshow[4.5mm], size=s).
% Related: polypeps_cell_holes (the first generation of tiles inside G, after the
% colouring drawn here); tn2d_two_family_purity (the two families). Neither
% draws the hypotheses of the lemma.
\begin{tenkzequation}
    \tenkzkernel
    \begin{tenkz}[size=s, rows={wire,ket,ket,ket,ket,ket,ket,ket,ket,wire}, cols=8,
            bonds=none]
        \tndeclare{species}{fam0}{hue=source:teal}
        \tn[at=(1,4), skin=none]{} \tn[at=(1,7), skin=none]{}
        \tn[at=(10,1), skin=none]{} \tn[at=(1,8), skin=none]{}
        \tnmark[form=enclosure, species=fam0, tint]
        {{(2,3), (2,4), (3,3), (3,4), (4,1), (4,2), (5,1), (5,2), (6,5), (6,6), (7,5), (7,6),
                    (8,1), (8,2), (9,1), (9,2)}}{}
        \tnmark[form=enclosure, species=secondary, tint]{(2,7) .. (5,8)}{}
        \tnmark[form=enclosure, species=selected, tint]{(4,3) .. (5,4)}{}
        \tnmark[form=label, species=selected, label pos=e]{midway (4,4) and (5,4)}{$G$}
        \tnmark[form=label, species=fam0, label pos=s]{midway (10,1) and (10,2)}{$B$}
        \tnmark[form=label, species=secondary, label pos=n]{midway (1,7) and (1,8)}{$C$}
        \tnmark[form=bracket, label pos=n]{(1,4) .. (1,7)}{$\ge cs$}
    \end{tenkz}
    $\longrightarrow$
    \begin{tenkz}[size=s, rows={wire,ket,ket,ket,ket,ket,ket,ket,ket,wire}, cols=8,
            bonds=none]
        \tndeclare{species}{fam0}{hue=source:teal}
        \tndeclare{species}{fam1}{hue=source:olive}
        \tndeclare{species}{centre}{hue=source:red}
        \tn[at=(1,4), skin=none]{} \tn[at=(10,4), skin=none]{}
        \tnmark[form=enclosure, species=fam1, tint]{(2,1) .. (9,8) - (4,3) .. (5,4)}{}
        \tnmark[form=enclosure, species=fam0, tint]
        {{(2,3), (2,4), (3,3), (3,4), (4,1), (4,2), (5,1), (5,2), (6,5), (6,6), (7,5), (7,6),
                    (8,1), (8,2), (9,1), (9,2)}}{}
        \tnmark[form=enclosure, species=secondary]{(2,7) .. (5,8)}{}
        \tn[at=0.5 n of 0.5 w of (4,3), skin=dot, species=centre, size=l]{}
        \tn[at=0.5 n of 0.5 e of (4,4), skin=dot, species=centre, size=l]{}
        \tn[at=0.5 s of 0.5 w of (5,3), skin=dot, species=centre, size=l]{}
        \tn[at=0.5 s of 0.5 e of (5,4), skin=dot, species=centre, size=l]{}
        \tnmark[form=label, species=fam0, label pos=n]{(1,4)}{$\text{forbidden }0$}
        \tnmark[form=label, species=fam1, label pos=s]{(10,4)}{$\text{forbidden }1$}
    \end{tenkz}
\end{tenkzequation}
